\documentclass{cv}

\usepackage[utf8]{inputenc}
\usepackage[ngerman]{babel}
\usepackage{hyperref}

\hypersetup{
	colorlinks=true,
	urlcolor=myblue,
	linkcolor=myblue,
	citecolor=myblue,
	pdftitle={Dr. André Kaderli -- Lebenslauf (Data \& AI)},
	pdfauthor={André Kaderli}
}

% Rote Platzhalter \tbd{...} vor dem Versand ausfuellen oder entfernen.

\begin{document}

	\name{Dr. André Kaderli}
	\tagline{Principal Data \& AI Engineer mit Promotion in theoretischer Physik und über fünf Jahren Erfahrung in End-to-End-Daten- und KI-Projekten für öffentliche Verwaltung und Privatwirtschaft -- von modernen Datenplattformen über prädiktive Modelle bis zu KI-Agenten. Ich verbinde mathematische Tiefe mit pragmatischer Umsetzung, führe interdisziplinäre Teams und berate Entscheidungsträger.}
	\begin{contactline}
		Obere Bahnhofstrasse 22M, 3700 Spiez\sep
		\href{mailto:andre.kaderli@gmail.com}{andre.kaderli@gmail.com}\sep
		\href{tel:+41793323931}{+41\,79 332 39 31}\sep
		\href{https://www.linkedin.com/in/andre-kaderli/}{linkedin.com/in/andre-kaderli}
	\end{contactline}

	\datedsection{Kernkompetenzen}{}
	\kv{Plattformen \& Engineering:}{Azure, Microsoft Fabric, Databricks, Apache Spark, Delta Lake, Power BI; Datenpipelines, Data Governance, CI/CD, Testing, Monitoring}
	\kv{Machine Learning \& KI:}{Prognosemodelle, Mustererkennung und Anomalieerkennung, KI-Agenten und LLM-Anwendungen}
	\kv{Programmierung:}{Python, SQL, PySpark, R, C++, MATLAB, Mathematica}
	\kv{Quantitative Methoden:}{Statistik, Optimierung, mathematische Modellierung und Simulation, Zeitreihenanalyse, Numerik}

	\datedsection{Berufserfahrung}{}

	\datedsubsection{Novalytica AG, Bern}{}

	\datedsubsubsection{Principal Data \& AI Engineer und Partner}{Seit Jan. 2026}
	\begin{itemize}
		\item Fachliche Leitung von Daten- und KI-Initiativen in öffentlicher Verwaltung und Privatwirtschaft; Führung und Mentoring interdisziplinärer Teams (\tbd{X} Personen)
		\item Konzeption, Aufbau und Betrieb moderner Daten- und Analytics-Plattformen (Microsoft, Databricks, Open Source) und Data-Governance-Frameworks
		\item Beratung von Entscheidungsträgern zu Datenstrategie und KI-Einsatz; \tbd{prüfen:} Mitverantwortung für Angebote, Kundenbeziehungen und Weiterentwicklung des Unternehmens als Partner
	\end{itemize}

	\datedsubsubsection{Senior Data Scientist}{Jan. 2023 -- Dez. 2025}
	\begin{itemize}
		\item Leitung von \tbd{X} End-to-End-Datenprojekten: Datenpipelines, BI-Reports, produktive ML-Modelle und intelligente Datenagenten
		\item Statistische Analysen grosser, heterogener Datensätze und mathematische Modellierung komplexer Systeme; \tbd{Ergebnis/Wirkung, z.\,B. Automatisierung, Zeitersparnis}
		\item Verantwortung für Entwicklungsstandards (CI/CD, Testing, Monitoring) im Team
	\end{itemize}

	\datedsubsubsection{Data Scientist}{Mai 2022 -- Dez. 2022}
	\begin{itemize}
		\item Data-Analytics- und BI-Projekte (Power BI, Azure, Databricks); Entwicklung von ML-Modellen und statistischen Auswertungen
	\end{itemize}

	\datedsubsection{Eraneos AG, Bern}{}
	\datedsubsubsection{Consultant}{Sep. 2021 -- Apr. 2022}
	\begin{itemize}
		\item IT- und Datenprojekte im öffentlichen Sektor mit Fokus auf Business-Analyse und quantitative Datenauswertungen
	\end{itemize}

	\datedsubsection{Max-Planck-Institut für Gravitationsphysik, Potsdam}{}
	\datedsubsubsection{Wissenschaftlicher Mitarbeiter (Doktorand)}{Sep. 2018 -- Aug. 2021}
	\begin{itemize}
		\item Forschung in theoretischer Physik: analytische und numerische Methoden, rekursive Algorithmen und Simulationsverfahren zur Lösung komplexer Differentialgleichungen; Veröffentlichungen (\href{https://inspirehep.net/authors/1809024}{inspirehep.net}) und Konferenzbeiträge
	\end{itemize}

	\newpage
	\datedsection{Zertifizierungen}{}
	\kv{Microsoft:}{Azure Data Scientist Associate, Azure Data Engineer Associate, Power BI Data Analyst Associate, Fabric Analytics Engineer Associate, Fabric Data Engineer Associate}
	\kv{Databricks:}{Databricks Data Engineer Professional, Databricks Data Engineer Associate}
	\kv{Projektmanagement:}{Hermes Advanced}

	\datedsection{Ausbildung}{}

	\datedsubsection{PhD in Physik, Humboldt-Universität zu Berlin}{Sep. 2018 -- Aug. 2021}
	\desc{Theoretische Physik (summa cum laude)}

	\datedsubsection{MSc in Physik, ETH Zürich}{Sep. 2016 -- Aug. 2018}
	\desc{Theoretische und Mathematische Physik (Abschlussnote: 6.0/6.0)}

	\datedsubsection{BSc in Physik, ETH Zürich}{Sep. 2013 -- Aug. 2016}
	\desc{Physik, Mathematik und Informatik (Abschlussnote: 5.7/6.0)}

	\datedsection{Wissensvermittlung}{}

	\datedsubsection{Data Academy, Novalytica}{Seit \tbd{Monat} 2022}
	\desc{Kurse zu Data Literacy, Programmierung (Python, R), Statistik, Datenanalyse und Visualisierung; Kundenschulungen zu eingeführten Datenplattformen (Microsoft Fabric, Power BI, Databricks)}

	\datedsubsection{Hochschullehre, HU Berlin und ETH Zürich}{2014 -- 2021}
	\desc{Übungsgruppen und Vorlesungsunterstützung in Linearer Algebra, Funktionentheorie und Quantenfeldtheorie; Betreuung von Studierenden und wissenschaftlichen Arbeiten}

	\datedsection{Weitere Angaben}{}
	\kv{Sprachen:}{Deutsch (Muttersprache), Englisch (verhandlungssicher), Französisch (Grundkenntnisse)}
	\kv{Stiftungen \& Stipendien:}{Alumni der Schweizerischen Studienstiftung, ehem. Träger des Ernst-Göhner-Stipendiums}
	\kv{Freiwilliges Engagement:}{Tourenleiter beim Schweizer Alpen-Club (SAC)}
	\kv{Interessen:}{Bergsteigen, Klettern, Skifahren, Gleitschirmfliegen, Jiu-Jitsu}

\end{document}
